\section{Global recovery of occupation and translation periods}
\label{sec:global-response}
The exact inverse admits a formulation without a periodic presentation.
Let $\mathcal G(D,d)$ consist of the nonempty, locally finite unions
$\mathcal O=\bigcup_{C\in\mathcal C}C$ of compact convex bodies with
nonempty interiors, diameters at most $D$, and distances between
distinct bodies at least $d>0$. The number of bodies and their shapes
need not be finite modulo translations. No boundary smoothness is
assumed in this section. Let $A$ be a compact convex body with nonempty
interior and diameter at most $\Delta$, and let $j$ be a probability
density supported on $A$ and positive almost everywhere in its
interior. The position of $A$ is unrestricted and its support function
need not be given. Fix the compass step so that
\begin{equation}\label{eq:global-step}
                         t+\Delta<d.
\end{equation}
Forward and reciprocal reverse preparations have the same meaning as
in Section~\ref{sec:setting}, with stationary displacement $Z\sim j$.
Write their exact mean fields as $F,R$, and put
\begin{equation}\label{eq:global-fields}
 v(x)=\int_A j(z)\one_{\mathcal O}(x+z)\dd z,\qquad
 g=F-R,\qquad
 Tf(x)=\frac14\sum_{e\in\{\pm e_1,\pm e_2\}}f(x+te).
\end{equation}
The short-segment identity gives $g=(T-I)v$. It uses no periodicity:
a segment of length $t<d$ cannot meet two different solids, and the
forward-minus-reverse bit is its endpoint occupation difference.
The positivity argument of Lemma~\ref{lem:noise-support} likewise gives
\begin{equation}\label{eq:global-positive-components}
 \{v>0\}=\bigcup_{C\in\mathcal C}\operatorname{int}P_C,
 \qquad P_C=C+(-A).
\end{equation}
These components have diameter at most $D+\Delta$ and mutual distance
at least $d-\Delta>t$. Also $0\le v\le1$. Translation continuity of
$j$ in $L^1$ implies uniform continuity of $v$, even when the family
of solids is infinite.

\begin{lemma}[Uniform exit from a positive component]
\label{lem:global-exit}
Let $X_k$ be the virtual compass walk with transition operator $T$,
started at $x\in\operatorname{int}P_C$, and let $\tau_C$ be its first
exit from that interior. Put
\begin{equation}\label{eq:global-exit-constants}
 H=\frac{(D+\Delta+t)^2}{t^2},\qquad b=\lceil2H\rceil.
\end{equation}
Then $v(X_{\tau_C})=0$ almost surely,
$\E_x\tau_C\le H$, and
\begin{equation}\label{eq:global-exit-tail}
                   \Prob_x(\tau_C>N)\le2^{-\lfloor N/b\rfloor}.
\end{equation}
The same estimates hold for every starting point and positive
component in the stated class. At a point where $v(x)=0$ set
$\tau_C=0$.
\end{lemma}
\begin{proof}
A step of length $t$ cannot pass from one positive component to
another. Hence the exit point has zero occupation, including when it
lies on the boundary. The process
$|X_k-x|^2-t^2k$ is a martingale. Before exit its distance from $x$ is
at most $D+\Delta$, and the exit step increases this bound by at most
$t$. Optional stopping at $n\wedge\tau_C$ therefore gives
$t^2\E_x(n\wedge\tau_C)\le(D+\Delta+t)^2$.
Monotone convergence proves the asserted mean bound and almost sure
exit. Markov's inequality bounds survival for $b$ steps by $1/2$.
At each surviving multiple of $b$, the same bound applies from the
current point in the same component. The Markov property proves
\eqref{eq:global-exit-tail}.
\end{proof}

\begin{theorem}[A global scalar-to-occupation inverse]
\label{thm:global-occupation}
Within the preceding class, $g$ determines $v$ uniquely, without
knowledge of the footprint or of the positive components. More
precisely, let $\mathcal T_H(x)$ be the stopping times of the virtual
compass walk, including zero, with $\E_x\tau\le H$. Then
\begin{equation}\label{eq:global-stopping-inverse}
 v(x)=\sup_{\tau\in\mathcal T_H(x)}
            \E_x\left[-\sum_{k=0}^{\tau-1}g(X_k)\right].
\end{equation}
In particular, two admissible occupation fields with forcing $g$ and
$\widetilde g$, possibly arising from different configurations,
footprints and densities satisfying the same bounds, obey
\begin{equation}\label{eq:global-inverse-stability}
                  \|v-\widetilde v\|_\infty
                         \le H\|g-\widetilde g\|_\infty.
\end{equation}
An explicit inverse is obtained from
\begin{equation}\label{eq:global-obstacle-iteration}
 v_0=0,\qquad v_{N+1}=\max\{0,Tv_N-g\}.
\end{equation}
It satisfies, on the whole plane,
\begin{equation}\label{eq:global-obstacle-error}
             0\le v-v_N\le2^{-\lfloor N/b\rfloor}.
\end{equation}
\end{theorem}
\begin{proof}
Since $(T-I)v=g$, the process
$v(X_n)-\sum_{k<n}g(X_k)$ is a martingale. For an integrable stopping
time, stopping first at $n\wedge\tau$ and then passing to the limit
is justified by boundedness of $v,g$ and integrability of $\tau$.
Thus
\[
 v(x)=\E_xv(X_\tau)
             -\E_x\sum_{k<\tau}g(X_k)
       \ge-\E_x\sum_{k<\tau}g(X_k).
\]
If $v(x)>0$, the exit time in Lemma~\ref{lem:global-exit} belongs to
$\mathcal T_H(x)$ and attains equality. If $v(x)=0$, the zero stopping
time does. This proves \eqref{eq:global-stopping-inverse}. The set of
stopping rules and the budget $H$ depend only on the stated bounds.
Each payoff changes by at most $H\|g-\widetilde g\|_\infty$ when
the forcing changes. Taking suprema proves
\eqref{eq:global-inverse-stability}, and also proves uniqueness.

Finite-horizon conditioning shows that \eqref{eq:global-obstacle-iteration}
is the value of the same payoff maximized over $\tau\le N$.
Consequently $0\le v_N\le v$. The particular rule
$\tau_C\wedge N$ gives
\[
 v(x)-v_N(x)\le\E_xv(X_{\tau_C\wedge N})
                     \le\Prob_x(\tau_C>N).
\]
The exit-tail estimate proves \eqref{eq:global-obstacle-error}.
This is the classical stopped Poisson argument, applied with the
unknown positive components discovered by the obstacle iteration;
the martingale and killed-walk identities are discussed in
\cite{LawlerLimic}.
\end{proof}

\begin{corollary}[Finite-field locality]
\label{cor:global-field-locality}
The value $v_N(x)$ depends only on the forcing at lattice points
reachable from $x$ in at most $N$ compass steps. If values
$\widetilde g$ at these points have absolute error at most $\epsilon$,
and $\widetilde v_N$ is computed by the same finite iteration, then
\begin{equation}\label{eq:global-field-locality}
 |v(x)-\widetilde v_N(x)|
             \le2^{-\lfloor N/b\rfloor}+N\epsilon.
\end{equation}
\end{corollary}
\begin{proof}
Reachability follows inductively from the transition rule. Averaging
and the positive-part map are nonexpansive in the supremum norm, so
the error in the $N$th iterate is at most $N\epsilon$. Combine this
with \eqref{eq:global-obstacle-error}.
\end{proof}

For a field or collection of fields $f$, let
$\operatorname{Per}(f)=\{w:f(x+w)=f(x)\text{ for all }x\}$, with
equality required in every component of a collection.

\begin{theorem}[Response rigidity for the full period group]
\label{thm:global-response-rigidity}
For every configuration and stationary law above,
\begin{equation}\label{eq:global-period-identity}
 \operatorname{Per}(F,R)=\operatorname{Per}(g)
      =\operatorname{Per}(v)=\operatorname{Per}(\mathcal O).
\end{equation}
Thus one exact reciprocal difference field determines the full
translation period group even when the footprint and its density are
unknown. This group is discrete. In particular, the configuration has
a rank-two lattice of periods, with finitely many component orbits,
if and only if the scalar difference field has two independent periods.
No periodicity or bounded-presentation hypothesis is needed for this
equivalence.
\end{theorem}
\begin{proof}
Every period of $\mathcal O$ is a period of both collision mean
fields, by translation covariance of the experiment. Every common
period of $F,R$ is a period of $g$. The inverse in
Theorem~\ref{thm:global-occupation} commutes with translations; hence
every period of $g$ is a period of $v$.

A period $w$ of $v$ preserves its positive set. The separated closures
of the positive components are precisely the bodies $P_C$, so
translation by $w$ permutes them. If $P_C+w=P_{C'}$, equality of
laboratory support functions gives, for every unit vector $u$,
\[
 p_C^{\rm lab}(u)+p_{-A}^{\rm lab}(u)+w\cdot u
       =p_{C'}^{\rm lab}(u)+p_{-A}^{\rm lab}(u).
\]
Cancellation yields $C+w=C'$. Applying the inverse translation proves
that $w$ is a period of $\mathcal O$, completing the equality chain.

A nonzero period cannot fix one compact component under translation.
It therefore takes that component to a distinct component, and its
length is at least $d$. The period group is consequently discrete.
If it contains two independent vectors, it is a rank-two lattice;
local finiteness and bounded component diameters imply that only
finitely many component orbits meet a compact fundamental cell.
Conversely such a periodic configuration has two independent periods,
which are inherited by $g$ through \eqref{eq:global-period-identity}.
\end{proof}

The exact datum here is a scalar response field on the plane. The
locality estimate specifies its finite-region approximation; it does
not replace a statistical count of the attempted bits used to estimate
those field values. Uniform finite period decisions from a fixed
protected aperture are a separate conclusion and retain the bounded
presentation and positive patch margin specified in
Section~\ref{sec:setting}. The present theorem identifies crystallinity
from the exact active response without assuming crystallinity of the
underlying configuration.
